% LaTeX master Datei(en) zusammengestellt von Falk-Jonatan Strube zur Nutzung an der Hochschule für Technik und Wirtschaft Dresden: https://github.com/genericFJS/htw
\documentclass{scrreprt}
\gdef\pathtomaster{../_LaTeX_master}
%% HTW Dresden Corporate Design
%% Copyright (C) Falk-Jonatan Strube <falkjonatanstrube@gmail.com>, 2017
%%% === Anfang Paket htwscr ===
\makeatletter
\ifcsname KOMAoption\endcsname%
\KOMAoption{footheight}{11mm}	% mehr Platz für Wasserzeichen o.ä. in der Fußzeile
\recalctypearea	% s.o.
\RequirePackage{geometry}
\geometry{top=25mm, bottom=25mm, headsep=10mm, inner=25mm, outer=25mm, footskip=15mm, marginparsep=2mm}
\RequirePackage{helvet}	% Corporate Design schreibt Arial/Helvetica vor
\renewcommand{\familydefault}{\sfdefault}	% s.o.
\RequirePackage{graphicx}	% für Grafiken
\RequirePackage{xcolor}		% für Farben
\definecolor{htworange}{RGB}{249,155,28}	% Corporate Design orange CMYK(0/45/100/0)
\definecolor{htwblue}{RGB}{0,116,188}			% Corporate Design blau CMYK(90/50/0/0)
\definecolor{htwgrey}{RGB}{128,128,128}		% Corporate Design gray 60% Schwarz [Annahme: CMYK(0,0,0,60)]
\RequirePackage{ifthen,xifthen}	% für Bedingungen/Optionen
\RequirePackage{etoolbox}				% s.o.

%%% Format des HTW-Headers (wird hier als Texteelement verwenden) nach dem Style der PowerPoint-Vorlagen. Weiterhin inspiriert durch das Paket tudscr (Corporate Design der TU Dresden)
%%% Dateipfad des Logos abhängig von der Verwendung als Paket/Input
\ifdefined\@isHTWCDpackage%
	\gdef\@pathtoimage{HTW}%
\else%
	\ifdefined\pathtomaster%
		\gdef\@pathtoimage{\pathtomaster /HTW}%
	\else 
		\gdef\@pathtoimage{HTW}%
	\fi%
\fi%
\newcommand{\pathtologo}[1]{\gdef\@pathtoimage{#1}}
%%% Berechnung der Größe des Logos
\newlength{\htwLogoHeight}
\setlength{\htwLogoHeight}{12mm}	% Mindesthöhe des Logos laut Corporate Design: 11mm
\newlength{\@htwLogoWidth}
\AtBeginDocument{
\def\printLogo{\includegraphics[height=\htwLogoHeight]{\@pathtoimage}}
\setlength\@htwLogoWidth{\widthof{\printLogo}}
}
%%% Definition des Header
\newcommand{\@htwheader}{{\noindent%
%%% Header ist abhängig von dem Vorhandensein von department, institute und chair
\newcommand{\hasDepInsCh}{\gdef\@hasDepInsCh{}}%
\ifdefined\@department\hasDepInsCh\fi%
\ifdefined\@institute\hasDepInsCh\fi%
\ifdefined\@chair\hasDepInsCh\fi%
\begin{minipage}[b][\htwLogoHeight]{\dimexpr \textwidth - \@htwLogoWidth - 2mm \relax}% 2mm Abstand zwischen Logo und Linie
{\small{\bfseries%
\ifdefined\@faculty\@faculty{}\fi} %
\strut}%
{\color{htworange}\hrule height 1pt}%
\vspace{.1mm}%
\end{minipage}
\hfill%
\begin{minipage}[b][\htwLogoHeight]{\@htwLogoWidth}
\includegraphics[height=\htwLogoHeight]{\@pathtoimage}
\end{minipage}
\ifdefined\@hasDepInsCh%
\\%
\begin{minipage}{\textwidth}
\small%
\ifdefined\@department\@department{}\fi%
\ifdefined\@institute\ifdefined\@department,~\fi \@institute\fi%
\ifdefined\@chair\ifdefined\@department,~\else\ifdefined\@institute,~\fi\fi\@chair\fi\strut
%{\color{htworange}\hrule height 0.5pt}
\end{minipage}
\fi
}}

%%% Definition der Fußzeilen
\usepackage{scrlayer-scrpage}
\@ifclassloaded{scrreprt}{%
\renewcommand*{\partpagestyle}{scrheadings}
\renewcommand*{\chapterpagestyle}{scrheadings}%
}{}
\pagestyle{scrheadings}
\cfoot{\vfill\pagemark}
\chead{}

%%% Einstellungen
\newcommand{\simpledate}{\gdef\@simpledate{}}		% Datum nicht mit "Eingereicht am:" Vorsatz
\newcommand{\nouppercase}{\gdef\@nouppercase{}}	% Titel/Parts/Chapters usw. nicht in Großbuchstaben

%%% Einiges erst zu Beginn des Dokumentes ausführen, um Einstellungen zu berücksichtigen
\AtBeginDocument{
\renewcommand\SS{SS}	% fix big ß
%%% Sektionen usw. ggf. in Großbuchstaben:
\providecommand{\chapterlinesformat}[3]{}	% compatibility
\providecommand{\chapterlineswithprefixformat}[3]{}	% compatibility
\providecommand{\sectionlinesformat}[4]{}	% compatibility
\ifdefined\@nouppercase\relax\else
\renewcommand\sectionlinesformat[4]{\@hangfrom{\hskip#2 #3}{\MakeUppercase{#4}}}
\@ifclassloaded{scrreprt}{%
\renewcommand\chapterlinesformat[3]{\@hangfrom{#2}{\MakeUppercase{#3}}}
\renewcommand\chapterlineswithprefixformat[3]{\MakeUppercase{#2#3}}%
}{}
\fi
%%% Anpassung der Part-Seite
\renewcommand*{\partheadstartvskip}{\@htwheader\vspace*{30mm}}
\renewcommand*{\raggedpart}{\raggedright}
\renewcommand*{\partformat}{\huge\ifdefined\@nouppercase \partname~\thepart\else\MakeUppercase{\partname~\thepart}\fi.}
\ifdefined\@nouppercase\setkomafont{part}{\Huge}\else\setkomafont{part}{\Huge\MakeUppercase}\fi
}

\newcommand{\bilingual}[3]{\ifdefined #1 \let #1\undefined\fi\newcommand{#1}{\iflanguage{english}{#3}{#2}}}	% Übersetzung von einigen Variablen in mehrere Sprachen

\newcommand{\checkThesis}[2]{\ifthenelse{\equal{#1}{diss}\OR\equal{#1}{doctoral}\OR\equal{#1}{phd}}{
\gdef#2{\dissertationname}}{
\ifthenelse{\equal{#1}{diploma}}{
\gdef#2{\diplomathesisname}}{
\ifthenelse{\equal{#1}{master}}{
\gdef#2{\masterthesisname}}{
\ifthenelse{\equal{#1}{bachelor}}{
\gdef#2{\bachelorthesisname}}{
\ifthenelse{\equal{#1}{student}}{
\gdef#2{\studentthesisname}}{
\ifthenelse{\equal{#1}{evidence}}{
\gdef#2{\studentresearchname}}{
\ifthenelse{\equal{#1}{project}}{
\gdef#2{\projectpapername}}{
\ifthenelse{\equal{#1}{seminar}}{
\gdef#2{\seminarpapername}}{
\ifthenelse{\equal{#1}{term}}{
\gdef#2{\termpapername}}{
\ifthenelse{\equal{#1}{research}}{
\gdef#2{\researchname}}{
\ifthenelse{\equal{#1}{log}}{
\gdef#2{\logname}}{
\ifthenelse{\equal{#1}{report}}{
\gdef#2{\reportname}}{
\ifthenelse{\equal{#1}{internship}}{
\gdef#2{\internshipname}}{
\ifthenelse{\equal{#1}{lecture}\OR\equal{#1}{lesson}\OR\equal{#1}{pract}}{
\bilingual{\professorname}{Vorlesung von}{Lecture by}
\simpledate
\gdef\@islecture{\relax}
\ifthenelse{\equal{#1}{lecture}}{\gdef#2{\lecturename}}{
\ifthenelse{\equal{#1}{lesson}}{\gdef#2{\lessonname}}{
\gdef#2{\practname}
}}}{
\gdef#2{#1}}}}}}}}}}}}}}}}

%%% Textfragmente der Titelseite in deutsch und englisch
\bilingual{\dateofbirthname}{Geboren am:}{Born on:}
\bilingual{\placeofbirthname}{in}{in}
\bilingual{\coursename}{Studiengang:}{Course:}
\bilingual{\disciplinename}{Studienrichtung:}{Discipline:}
\bilingual{\matriculationnumbername}{Matrikelnummer:}{Matriculation number:}
\bilingual{\matriculationyearname}{Immatrikulationsjahr}{Matriculation year:}
\bilingual{\graduationname}{zur Erlangung des akademischen Grades}{to achieve the academic degree}
\bilingual{\refereename}{Gutachter}{Referee}
\bilingual{\advisorname}{Fachreferent}{Advisor}
\bilingual{\supervisorname}{Betreuer}{Supervisor}
\bilingual{\professorname}{Betreuender Hochschullehrer}{Supervising professor}
\bilingual{\datename}{Eingereicht am:}{Submitted on:}
%% Thesis Typen (inspiriert von dem Paket tudscr):
\bilingual{\dissertationname}{Dissertation}{Dissertation} % diss, doctoral, phd 
\bilingual{\diplomathesisname}{Diplomarbeit}{Diploma Thesis} % diploma 
\bilingual{\masterthesisname}{Master-Arbeit}{Master Thesis} % master 
\bilingual{\bachelorthesisname}{Bachelor-Arbeit}{Bachelor Thesis} % bachelor 
\bilingual{\studentthesisname}{Studienarbeit}{Student Thesis} % student 
\bilingual{\studentresearchname}{Großer Beleg}{Student Research Project} % evidence 
\bilingual{\projectpapername}{Projektarbeit}{Project Paper} % project
\bilingual{\seminarpapername}{Seminararbeit}{Seminar Paper} % seminar
\bilingual{\termpapername}{Hausarbeit}{Term Paper} % term
\bilingual{\researchname}{Forschungsbericht}{Research Report} % research 
\bilingual{\logname}{Protokoll}{Log} % log 
\bilingual{\reportname}{Bericht}{Report} % report
\bilingual{\internshipname}{Praktikumsbericht}{Internship Report} % internship 
\bilingual{\lecturename}{Vorlesungsmitschrift}{Lecture Notes} % lecture 
\bilingual{\lectureauthorname}{Mitschrift von}{Notes by} % lecture  
\bilingual{\lessonname}{\"Ubungsmitschrift}{Lesson Notes} % lecture  
\bilingual{\practname}{Praktikumsmitschrift}{Pracital Course Notes} % lecture 
%%% Variablen für die Titelseite (inspiriert von dem Paket tudscr):
\renewcommand{\subject}[1]{\checkThesis{#1}{\@subject}}
\newcommand{\faculty}[1]{\gdef\@faculty{#1}}
\newcommand{\department}[1]{\gdef\@department{#1}}
\newcommand{\institute}[1]{\gdef\@institute{#1}}
\newcommand{\chair}[1]{\gdef\@chair{#1}}
\newcommand{\authormore}[1]{\gdef\@authormore{#1}}
\newcommand{\dateofbirth}[1]{\gdef\@dateofbirth{#1}}
\newcommand{\placeofbirth}[1]{\gdef\@placeofbirth{#1}}
\newcommand{\course}[1]{\gdef\@course{#1}}
\newcommand{\discipline}[1]{\gdef\@discipline{#1}}
\newcommand{\matriculationnumber}[1]{\gdef\@matriculationnumber{#1}}
\newcommand{\matriculationyear}[1]{\gdef\@matriculationyear{#1}}
\newcommand{\thesis}[1]{\checkThesis{#1}{\@thesis}}
\newcommand{\graduation}[2][]{\gdef\@graduation{#2}\gdef\@graduationshort{#1}}
\newcommand{\supervisor}[1]{\gdef\@supervisor{#1}}
\newcommand{\referee}[1]{\gdef\@referee{#1}}
\newcommand{\advisor}[1]{\gdef\@advisor{#1}}
\newcommand{\professor}[1]{\gdef\@professor{#1}}
%%% Titelseite (inspiriert von dem Paket tudscr):
\renewcommand*{\maketitle}{%
\begin{titlepage}%
\begin{flushleft}
\renewcommand{\thefootnote}{$\star$}%
%%% Kopfzeile HTW-Style:
\@htwheader\\[13mm]
%%% Die Zeilen für titlehead, subject, title, subtitle und author sind immer vorhanden (spacing)
{\noindent\@titlehead{}\strut}\\[3mm]%
{\Large\@subject{}\strut}\\[8mm]%
{\Huge\bfseries%
\ifdefined\@nouppercase\@title{}\else\MakeUppercase{\@title{}}\fi%
\strut}\\
{\Large\bfseries%
\ifdefined\@nouppercase\@subtitle{}\else\MakeUppercase{\@subtitle{}}\fi%
\strut}\\[11mm]
\ifdefined\@islecture\strut\else{\Large\@author{}\strut}\fi % Autor wird bei Vorlesungsmitschrift unten angezeigt
%%% Weitere Felder werden nur angezeigt, wenn sie definiert wurden
\ifdefined\@authormore\ifdefined\@islecture\else\\\fi{\@authormore{}\strut}\fi
\ifdefined\@dateofbirth\\\dateofbirthname{} \@dateofbirth{} \ifdefined\@placeofbirth\placeofbirthname{} \@placeofbirth{}\strut\fi\fi
\ifdefined\@course\\\coursename{} \@course{}\strut\fi
\ifdefined\@discipline\\\disciplinename{} \@discipline{}\strut\fi
\ifdefined\@matriculationnumber\\\matriculationnumbername{} \@matriculationnumber{}\strut\fi
\ifdefined\@matriculationyear\\\matriculationyearname{} \@matriculationyear{}\strut\fi
\ifdefined\@thesis\\[10mm]{\LARGE\bfseries
\ifdefined\@nouppercase\@thesis{}\else\MakeUppercase{\@thesis{}}\fi
\strut}
\ifdefined\@graduation\\[2mm]{\graduationname{}\strut}
\\[5mm]{\LARGE\bfseries
\ifdefined\@nouppercase\relax\else\MakeUppercase{\fi%
\@graduation{} \ifdefined\@graduationshort(\@graduationshort{})\fi
\ifdefined\@nouppercase\relax\else}\fi\strut}\fi\fi
\vfill\noindent
\ifdefined\@islecture\\[1mm]{\lectureauthorname{}\strut}\\[.7mm]{\Large\@author{}\strut}\fi
\ifdefined\@referee{\refereename{}\strut}\\[.7mm]{\Large\@referee{}\strut}\fi
\ifdefined\@advisor\\[1mm]{\advisorname{}\strut}\\[.7mm]{\Large\@advisor{}\strut}\fi
\ifdefined\@supervisor\\[1mm]{\supervisorname{}\strut}\\[.7mm]{\Large\@supervisor{}\strut}\fi
\ifdefined\@professor\\[1mm]{\professorname{}\strut}\\[.7mm]{\Large\@professor{}\strut}\fi
\ifdefempty{\@date}{}{\ifdefined\@simpledate\\[6mm]\relax\else\\[3mm]\datename{}~\fi\@date{}\strut}
\ifdefempty{\@publishers}{}{\\[4.5mm]{\Large\@publishers{}\strut}}%
\renewcommand\footnoterule{}
\@thanks\global\let\@thanks\@empty
\end{flushleft}
\end{titlepage}
\setcounter{footnote}{0}
}
\else%
\PackageError{htwcd}{htwcd needs a KOMA class!}%
\fi%
\makeatother
%%% === Ende Paket htwscr ===
%-- Nummerierung bis Subsubsection für Report
\setcounter{secnumdepth}{3}	% zählt auch subsubsection
\setcounter{tocdepth}{3}	% Inhaltsverzeichnis bis in subsubsection

% LuaLaTeX Unterscheidung
\usepackage{ifluatex}

% Input inkl. Umlaute, Silbentrennung
\ifluatex
	\usepackage{fontspec}
\else
	\usepackage[utf8]{inputenc}
	\usepackage[T1]{fontenc}
\fi
\usepackage[english,ngerman]{babel}
\usepackage{csquotes}	% Anführungszeichen
\usepackage{amssymb}
\usepackage{marvosym}
\usepackage{eurosym}

% Style-Aufhübschung
\usepackage{soul, color}	% Kapitälchen, Unterstrichen, Durchgestrichen usw. im Text
%\usepackage{titleref}
\usepackage[breakwords, fit]{truncate}	% Abschneiden von Sätzen
\renewcommand{\TruncateMarker}{\,…}
\RequirePackage{scrlayer-scrpage}	% Kopf- und Fußzeilen
\usepackage{setspace} % Zeilenabstand anpassen

% Mathe usw.
\usepackage{amssymb}
\usepackage{amsthm}
\usepackage[fleqn,intlimits]{amsmath}	% fleqn: align-Umgebung rechtsbündig; intlimits: Integralgrenzen immer ober-/unterhalb
%\usepackage{mathtools} % u.a. schönere underbraces
\usepackage{xcolor}
\usepackage{esint}	% Schönere Integrale, \oiint vorhanden
\everymath=\expandafter{\the\everymath\displaystyle}	% Mathe Inhalte werden weniger verkleinert
\usepackage{wasysym}	% mehr Symbole, bspw \lightning
%\renewcommand{\int}{\int\limits}
%\usepackage{xfrac}	% mehr fracs: sfrac{}{}
\let\oldemptyset\emptyset	% schöneres emptyset
\let\emptyset\varnothing
%\RequirePackage{mathabx}	% mehr Symbole
\mathchardef\mhyphen="2D	% Hyphen in Math
\usepackage{mathtools}		% zum Beispiel {pmatrix}[r] für Formatierung der Matrix

% tikz usw.
\usepackage{tikz}
\usepackage{pgfplots}
\pgfplotsset{compat=1.11}	% Umgeht Fehlermeldung
\usetikzlibrary{graphs}
%\usetikzlibrary{through}	% ???
\usetikzlibrary{arrows}
\usetikzlibrary{arrows.meta}	% Pfeile verändern / vergrößern: \draw[-{>[scale=1.5]}] (-3,5) -> (-3,3);
\usetikzlibrary{automata,positioning} % Zeilenumbruch im Node node[align=center] {Text\\nächste Zeile} automata für Graphen
\usetikzlibrary{matrix}
\usetikzlibrary{patterns}	% Schraffierte Füllung
\usetikzlibrary{shapes.geometric}	% Polygon usw.
\tikzstyle{reverseclip}=[insert path={	% Inverser Clip \clip
	(current page.north east) --
	(current page.south east) --
	(current page.south west) --
	(current page.north west) --
	(current page.north east)}
% Nutzen: 
%\begin{tikzpicture}[remember picture]
%\begin{scope}
%\begin{pgfinterruptboundingbox}
%\draw [clip] DIE FLÄCHE, IN DER OBJEKT NICHT ERSCHEINEN SOLL [reverseclip];
%\end{pgfinterruptboundingbox}
%\draw DAS OBJEKT;
%\end{scope}
%\end{tikzpicture}
]	% Achtung: dafür muss doppelt kompliert werden!
\usepackage{graphpap}	% Grid für Graphen
\tikzset{every state/.style={inner sep=2pt, minimum size=2em}}
\usetikzlibrary{mindmap, backgrounds}
%\usepackage{tikz-uml}	% braucht Dateien: http://perso.ensta-paristech.fr/~kielbasi/tikzuml/

% Tabular
\usepackage{longtable}	% Große Tabellen über mehrere Seiten
\usepackage{multirow}	% Multirow/-column: \multirow{2[Anzahl der Zeilen]}{*[Format]}{Test[Inhalt]} oder \multicolumn{7[Anzahl der Reihen]}{|c|[Format]}{Test2[Inhalt]}
\usepackage{multicol}
\renewcommand{\arraystretch}{1.3} % Tabellenlinien nicht zu dicht
\usepackage{colortbl}
\arrayrulecolor{gray}	% heller Tabellenlinien
\usepackage{array}	% für folgende 3 Zeilen (für Spalten fester breite mit entsprechender Ausrichtung):
\newcolumntype{L}[1]{>{\raggedright\let\newline\\\arraybackslash\hspace{0pt}}m{\dimexpr#1\columnwidth-2\tabcolsep-1.5\arrayrulewidth}}
\newcolumntype{C}[1]{>{\centering\let\newline\\\arraybackslash\hspace{0pt}}m{\dimexpr#1\columnwidth-2\tabcolsep-1.5\arrayrulewidth}}
\newcolumntype{R}[1]{>{\raggedleft\let\newline\\\arraybackslash\hspace{0pt}}m{\dimexpr#1\columnwidth-2\tabcolsep-1.5\arrayrulewidth}}
\usepackage{booktabs}	% Kompatibilität für array und \hrule
\usepackage{caption}	% Um auch unbeschriftete Captions mit \caption* zu machen

% Nützliches
\usepackage{verbatim}	% u.a. zum auskommentieren via \begin{comment} \end{comment}
\usepackage{tabto}	% Tabs: /tab zum nächsten Tab oder /tabto{.5 \CurrentLineWidth} zur Stelle in der Linie
\NumTabs{6}	% Anzahl von Tabs pro Zeile zum springen
\usepackage{listings} % Source-Code mit Tabs
\usepackage{lstautogobble} 
\usepackage{enumitem}	% Anpassung der enumerates
%\setlist[enumerate,1]{label=(\arabic*)}	% global andere Enum-Items
\renewcommand{\labelitemiii}{$\scriptscriptstyle ^\blacklozenge$} % global andere 3. Item-Aufzählungszeichen
\usepackage{letltxmacro} % neue Definiton von Grundbefehlen
% Nutzen:
%\LetLtxMacro{\oldemph}{\emph}
%\renewcommand{\emph}[1]{\oldemph{#1}}
\RequirePackage{xpatch}	% ua. Konkatenieren von Strings/Variablen (etoolbox)
\usepackage{xstring}	% String Operationen
\usepackage{minibox}	% Minibox anstatt \fbox{} für Boxen mit (manuellen) Zeilenumbruch
\usepackage{framed}		% \begin{framed} statt minibox/fbox mit automatischem Zeilenumbruch
\usepackage{microtype}	% Verfeinerte Hyphenation
\usepackage{import}		
\usepackage{standalone}	% Zusammen mit import: Einfügen von separaten Latex-Dateien (die schon header haben)
\usepackage{datetime}	% eigenes Datum bspw. mit \newdate{date}{06}{09}{2012} und dann \date{\displaydate{date}}
\usepackage{chngcntr} % Counter verändern, bspw. Abbildungen \counterwithout{figure}{chapter} und Tabellen \counterwithout{table}{chapter}

% Einrichtung von lst
%\renewcommand*{\ttdefault}{qcr}
\lstset{
basicstyle=\ttfamily, 
%basicstyle=\footnotesize\sffamily\color{black},
%mathescape=true, 
%escapeinside=^^, 
autogobble, 
tabsize=2,
frame=single,
rulecolor=\color{lightgray},
numbers=left,
numbersep=5pt,
numberstyle=\tiny\color{gray},
commentstyle=\color{gray},
keywordstyle=\color{htwblue},
stringstyle=\color{htworange},
morecomment=[l][\color{magenta}]{\#}
showspaces=false,
showstringspaces=false,
breaklines=true,
literate=%
    {Ö}{{\"O}}1
    {Ä}{{\"A}}1
    {Ü}{{\"U}}1
    {ß}{{\ss}}1
    {ü}{{\"u}}1
    {ä}{{\"a}}1
    {ö}{{\"o}}1
    {„}{{\glqq}}1
    {“}{{\grqq}}1
    {~}{{\textasciitilde}}1
}
\usepackage{scrhack} % Fehler umgehen
\def\ContinueLineNumber{\lstset{firstnumber=last}} % vor lstlisting. Zum wechsel zum nicht-kontinuierlichen muss wieder \StartLineAt1 eingegeben werden
\def\StartLineAt#1{\lstset{firstnumber=#1}} % vor lstlisting \StartLineAt30 eingeben, um bei Zeile 30 zu starten
\let\numberLineAt\StartLineAt

% BibTeX
\usepackage[
backend=biber,
sortlocale=de_DE, 
bibencoding=utf8,
%bibencoding=ascii,
style=authoryear, 
citestyle=authoryear-ibid,
%doi=false,
%isbn=false,
%url=false
]{biblatex}	% BibTeX
%\bibliographystyle{acm}
\usepackage{makeidx}
%\makeglossary
%\makeindex
\usepackage{url}	% Darstellung von URLs

% Grafiken
\usepackage{graphicx}
\usepackage{subcaption}	% subfigures (mehrere Graphiken in einer Figure
\usepackage{epstopdf}	% eps-Vektorgrafiken einfügen
\usepackage{transparent}	% transparent nutzen: {\transparent{0.4} ...}
%\epstopdfsetup{outdir=\customDir}
% Prüft, ob Grafik existiert (mit \ifvalidimage{}{}) [Quelle: https://tex.stackexchange.com/a/99176]:
\makeatletter
\newif\ifgraphicexist
\catcode`\*=11
\newcommand\ifvalidimage[1]{%
    \begingroup
    \global\graphicexisttrue
    \let\input@path\Ginput@path
    \filename@parse{#1}%
    \ifx\filename@ext\relax
    \@for\Gin@temp:=\Gin@extensions\do{%
        \ifx\Gin@ext\relax
        \Gin@getbase\Gin@temp
        \fi}%
    \else
    \Gin@getbase{\Gin@sepdefault\filename@ext}%
    \ifx\Gin@ext\relax
    \global\graphicexistfalse
    \def\Gin@base{\filename@area\filename@base}%
    \edef\Gin@ext{\Gin@sepdefault\filename@ext}%
    \fi
    \fi
    \ifx\Gin@ext\relax
    \global\graphicexistfalse
    \else 
    \@ifundefined{Gin@rule@\Gin@ext}%
    {\global\graphicexistfalse}%
    {}%
    \fi  
    \ifx\Gin@ext\relax 
    \gdef\imageextension{unknown}%
    \else
    \xdef\imageextension{\Gin@ext}%
    \fi 
    \endgroup 
    \ifgraphicexist
    \expandafter \@firstoftwo
    \else
    \expandafter \@secondoftwo
    \fi 
} 
\catcode`\*=12
\makeatother
\LetLtxMacro{\forceincludegraphics}{\includegraphics}	% neuer Befehl für includegraphics
\renewcommand{\includegraphics}[2][]{	% altes includegraphics neu definieren, damit es auch nicht vorhandene einfügt
\ifvalidimage{#2}{
\forceincludegraphics[#1]{#2}
}{
%\message{Achtung: Grafik wurde nicht gefunden: '#2'}
\ClassWarning{graphics}{Missing included graphic: '#2'}
\minibox[frame]{
\textbf{\StrSubstitute{#2}{_}{\_}}%  \ifthenelse{\isempty{#1}}{}{\\\textit{#1}}
}}}

% pdf-Setup
\usepackage{pdfpages}
\usepackage[bookmarks,%
bookmarksopen=false,% Klappt die Bookmarks in Acrobat aus
colorlinks=true,%
linkcolor=black,%
citecolor=black,%
urlcolor=htworange,%
]{hyperref}

% Nach Dokumentenbeginn ausführen:
\makeatletter
\AtBeginDocument{	
	% pdf Informationen
	\hypersetup{
		pdfauthor={\@author},
		pdftitle={\@title},
		pdfproducer={LaTeX Layout erstellt von Falk-Jonatan Strube}
	}
	% signierte Fußzeile
	%\providecommand{\ofootname}{Mitschrift von}
	\providecommand{\pathtomaster}{../}
	\cfoot{%
	%\ifdefined\ofoottext
	%\else
	%\begin{minipage}{40mm}
	%\begin{flushright}
	%\textcolor{gray}{
	%\ofootname\\
	%Falk-Jonatan Strube\vspace*{-1mm}}
	%\end{flushright}
	%\end{minipage}
	%\fi%
	%\hspace*{1mm}%
	\vspace*{-4mm}
	\begin{minipage}{10 mm}
	\vspace*{-2mm}
	\includegraphics[height=9mm]{\pathtomaster/FJS-logo_gray}
	\end{minipage}%
	}
}
\makeatother



%%% ==================================================
%%% ===================== MACROS =====================
%%% ==================================================
% Horizontale Linie:
\newcommand{\HRule}[1][\medskipamount]{\par
  \vspace*{\dimexpr-\parskip-\baselineskip+#1}
  \noindent\rule[0.2ex]{\linewidth}{0.2mm}\par
  \vspace*{\dimexpr-\parskip-.5\baselineskip+#1}}
% Gestrichelte horizontale Linie:
\RequirePackage{dashrule}
\newcommand{\HDRule}[1][\medskipamount]{\par
  \vspace*{\dimexpr-\parskip-\baselineskip+#1}
  \noindent\hdashrule[0.2ex]{\linewidth}{0.2mm}{1mm} \par
  \vspace*{\dimexpr-\parskip-.5\baselineskip+#1}}
% Mathe in Anführungszeichen:
\newsavebox{\mathbox}\newsavebox{\mathquote}
\makeatletter
\newcommand{\mq}[1]{% \mathquotes{<stuff>}
  \savebox{\mathquote}{\text{"}}% Save quotes
  \savebox{\mathbox}{$\displaystyle #1$}% Save <stuff>
  \raisebox{\dimexpr\ht\mathbox-\ht\mathquote\relax}{"}#1\raisebox{\dimexpr\ht\mathbox-\ht\mathquote\relax}{''}
}
\makeatother

% Paragraph mit Zähler (Section-Weise)
\newcounter{cparagraphC}
\newcommand{\cparagraph}[1]{
\stepcounter{cparagraphC}
\paragraph{\thesection{}-\thecparagraphC{} #1}
%\addcontentsline{toc}{subsubsection}{\thesection{}-\thecparagraphC{} #1}
\label{\thesection-\thecparagraphC}
}
\makeatletter
\@addtoreset{cparagraphC}{section}
\makeatother


% (Vorlesungs-)Folien einbinden:
\makeatletter
% Folien-Pfad (inkl. Dateiname ohne Endung und ggf. ohne Nummerierung)
\newcommand{\@slidePath}{}
\newcommand{\slidePath}[1]{\renewcommand{\@slidePath}{#1}}
% Folien Eigenschaften
\newcommand{\@slideScale}{0.5}
\newcommand{\slideScale}[1]{\renewcommand{\@slideScale}{#1}}
\AtBeginDocument{
% Folien von einer Datei skaliert
\newcommand{\slide}[2][\@slideScale]{\slides[#1]{}{#2}}
\newcommand{\slideTrim}[6][\@slideScale]{\slides[#1 , clip,  trim = #5cm #4cm #6cm #3cm]{}{#2}}
% Folien von mehreren nummerierten Dateien skaliert
\newcommand{\slides}[3][\@slideScale]{\begin{center}
\IfFileExists{\@slidePath #2.pdf}{
	\includegraphics[page=#3, scale=#1]{\@slidePath #2.pdf}
}{
	\ClassWarning{graphics}{Missing included slide: '\@slidePath #2.pdf'}
	\minibox[frame]{
		\textbf{\StrSubstitute{\@slidePath #2.pdf}{_}{\_}}\\
		Folie #3
	}
}
\end{center}}
}
\makeatother

% \emph{} anders definieren
\makeatletter
\DeclareRobustCommand{\em}{%
  \@nomath\em \if b\expandafter\@car\f@series\@nil
  \normalfont \else \scshape \fi}
\makeatother

% unwichtiges
\newcommand{\unimptnt}[1]{{\transparent{0.5}#1}}

% alph. enumerate
\newenvironment{anumerate}{\begin{enumerate}[label=(\alph*)]}{\end{enumerate}} % Alphabetische Aufzählung

% Hanging parameters
\newcommand{\hangpara}[1]{\par\noindent\hangindent+2em\hangafter=1 #1\par\noindent}

% Nicht-nummerierte Überschriften im Inhaltsverzeichnis
\newcommand{\chapterN}[1]{\chapter*{#1}%
\addcontentsline{toc}{chapter}{\protect\numberline{}#1}}
\newcommand{\sectionN}[1]{\section*{#1}%
\addcontentsline{toc}{section}{\protect\numberline{}#1}}
\newcommand{\subsectionN}[1]{\subsection*{#1}%
\addcontentsline{toc}{subsection}{\protect\numberline{}#1}}
\newcommand{\subsubsectionN}[1]{\subsubsection*{#1}%
\addcontentsline{toc}{subsubsection}{\protect\numberline{}#1}}

%% EINFACHE BEFEHLE

% Abkürzungen Mathe
\newcommand{\EE}{\mathbb{E}}
\newcommand{\QQ}{\mathbb{Q}}
\newcommand{\RR}{\mathbb{R}}
\newcommand{\CC}{\mathbb{C}}
\newcommand{\NN}{\mathbb{N}}
\newcommand{\ZZ}{\mathbb{Z}}
\newcommand{\PP}{\mathbb{P}}
\renewcommand{\SS}{\mathbb{S}}
\newcommand{\cA}{\mathcal{A}}
\newcommand{\cB}{\mathcal{B}}
\newcommand{\cC}{\mathcal{C}}
\newcommand{\cD}{\mathcal{D}}
\newcommand{\cE}{\mathcal{E}}
\newcommand{\cF}{\mathcal{F}}
\newcommand{\cG}{\mathcal{G}}
\newcommand{\cH}{\mathcal{H}}
\newcommand{\cI}{\mathcal{I}}
\newcommand{\cJ}{\mathcal{J}}
\newcommand{\cM}{\mathcal{M}}
\newcommand{\cN}{\mathcal{N}}
\newcommand{\cP}{\mathcal{P}}
\newcommand{\cR}{\mathcal{R}}
\newcommand{\cS}{\mathcal{S}}
\newcommand{\cZ}{\mathcal{Z}}
\newcommand{\cL}{\mathcal{L}}
\newcommand{\cT}{\mathcal{T}}
\newcommand{\cU}{\mathcal{U}}
\newcommand{\cX}{\mathcal{X}}
\newcommand{\cV}{\mathcal{V}}
\renewcommand{\phi}{\varphi}
\renewcommand{\epsilon}{\varepsilon}
\renewcommand{\theta}{\vartheta}

% Verschiedene als Mathe-Operatoren
\DeclareMathOperator{\arccot}{arccot}
\DeclareMathOperator{\arccosh}{arccosh}
\DeclareMathOperator{\arcsinh}{arcsinh}
\DeclareMathOperator{\arctanh}{arctanh}
\DeclareMathOperator{\arccoth}{arccoth} 
\DeclareMathOperator{\var}{Var} % Varianz 
\DeclareMathOperator{\cov}{Cov} % Co-Varianz 

% Farbdefinitionen
\definecolor{red}{RGB}{180,0,0}
\definecolor{green}{RGB}{75,160,0}
\definecolor{blue}{RGB}{0,75,200}
\definecolor{orange}{RGB}{255,128,0}
\definecolor{yellow}{RGB}{255,245,0}
\definecolor{purple}{RGB}{75,0,160}
\definecolor{cyan}{RGB}{0,160,160}
\definecolor{brown}{RGB}{120,60,10}

\definecolor{itteny}{RGB}{244,229,0}
\definecolor{ittenyo}{RGB}{253,198,11}
\definecolor{itteno}{RGB}{241,142,28}
\definecolor{ittenor}{RGB}{234,98,31}
\definecolor{ittenr}{RGB}{227,35,34}
\definecolor{ittenrp}{RGB}{196,3,125}
\definecolor{ittenp}{RGB}{109,57,139}
\definecolor{ittenpb}{RGB}{68,78,153}
\definecolor{ittenb}{RGB}{42,113,176}
\definecolor{ittenbg}{RGB}{6,150,187}
\definecolor{itteng}{RGB}{0,142,91}
\definecolor{ittengy}{RGB}{140,187,38}

% Textfarbe ändern
\newcommand{\tred}[1]{\textcolor{red}{#1}}
\newcommand{\tgreen}[1]{\textcolor{green}{#1}}
\newcommand{\tblue}[1]{\textcolor{blue}{#1}}
\newcommand{\torange}[1]{\textcolor{orange}{#1}}
\newcommand{\tyellow}[1]{\textcolor{yellow}{#1}}
\newcommand{\tpurple}[1]{\textcolor{purple}{#1}}
\newcommand{\tcyan}[1]{\textcolor{cyan}{#1}}
\newcommand{\tbrown}[1]{\textcolor{brown}{#1}}

% Umstellen der Tabellen Definition
\newcommand{\mpb}[1][.3]{\begin{minipage}{#1\textwidth}\vspace*{3pt}}
\newcommand{\mpe}{\vspace*{3pt}\end{minipage}}

\newcommand{\resultul}[1]{\underline{\underline{#1}}}
\newcommand{\parskp}{$ $\\}	% new line after paragraph
\newcommand{\parnl}{\unskip\parfillskip 0pt \par} % justified new line (brace line!)

% mathematische Abkürzungen
\newcommand{\corr}{\,\widehat{=}\,}
\newcommand{\mdeg}{^{\circ}}
\newcommand{\nok}[2]{\binom{#1}{#2}}	% n über k BESSER: \binom{n}{k}
\newcommand{\mtr}[1]{\begin{pmatrix}#1\end{pmatrix}}	% Matrix
\newcommand{\dtr}[1]{\begin{vmatrix}#1\end{vmatrix}}	% Determinante (Betragsmatrix)
\LetLtxMacro{\originalVec}{\vec}
\renewcommand{\vec}[1]{\underline{#1}}	% Vektorschreibweise
\newcommand{\imptnt}[1]{\colorbox{red!30}{#1}}	% Wichtiges
\newcommand{\intd}[1]{\,\mathrm{d}#1}
\newcommand{\diffd}[1]{\mathrm{d}#1}
% für Module-Rechnung: \pmod{}
\newcommand{\unit}[1]{\,\mathrm{#1}}
%\LetLtxMacro{\ntilde}{\tilde}
%\renewcommand{\tilde}{\widetilde}
\newcommand{\To}{\Rightarrow}

% weiter Abkürzungen
\newcommand{\gdw}{genau dann wenn}
\newcommand{\lecdate}[1]{\marginpar{\small\textcolor{gray}{Vorlesung\newline #1}}}
\newcommand{\lesdate}[1]{\marginpar{\small\textcolor{gray}{Übung\newline #1}}}

\faculty{Fakultät Informatik/Mathematik}
\chair{Lehrstuhl für Medientechnik / Audio- und Videosysteme}
\subject{lecture}
\title{Rechnerarchitektur}
\author{Falk-Jonatan Strube}
\professor{Prof. Dr. Jens Schönthier}

\makeatletter
\@addtoreset{chapter}{part}
\makeatother
\slidePath{Vorlesung/RA_043_}
\slideScale{0.7}

\begin{document}
\maketitle
\tableofcontents

%\chapter*{Vorbemerkung}
%Passwort Folien: \\% mi043##

% Folien i: eher weniger relevante Prüfungsinformationen
\part{Vorlesung}

\chapter{Architektur nach von Neumann}

\unimptnt{
%\slidesScale{01}{2}
\section*{Streiflicht: Konrad Zuse}
\slidesScale{01}{3}
\slidesScale{01}{4}
\section*{Streiflicht: John von Neumann}
\slidesScale{01}{5}

\section{Grundidee von von Neumann}
\slidesScale{01}{6}
}
\section{Grundarchitekturen von Rechnern/Computern}
\slidesScale{01}{7}
Havard-Architektur: Hardware-Aufwand größer durch mehr Busse

\section{Von-Neumann-Architektur}
\slidesScale{01}{8}
Im Speicher sind auf den ersten Blick von außen Daten und Befehle nicht voneinander zu unterschieden.
\slidesScale{01}{9}

\subsection{Kurzfassung Von-Neumann-Rechner Prinzip}
\slidesScale{01}{10}
\paragraph{Blockschaltbild}
\slidesScale{01}{11}
% Folie 12: Bild Wikipedia

\section{Befehlssequenz}
\slidesScale{01}{13}
Für Sprungbefehle werden für performante Ausführung Sprungvorhersagen benötigt.

\section{Operandenadressierung}
\slidesScale{01}{14}

\section{Architekturverfeinerung 1: Befehlszähler}
\slidesScale{01}{15}
\unimptnt{
$\Rightarrow$ Von-Neumann Rechner mit Befehlszähler
\slidesScale{01}{16}
}

\section{Klassischer Universalspeicher}
\subsection{Speicherwerk}
\slidesScale{01}{17}
\subsection{Rechenwerk}
\slidesScale{01}{18}
\unimptnt{
\subsubsection{Blockschaltbild einer ALU}
\slidesScale{01}{19}
}
\subsubsection{Flags/Bits des Statusregisters}
\slidesScale{01}{20}
$\mu P$: Mikroprozessor
\unimptnt{
\subsubsection*{Bsp.: Erzeugung des Bedingungsbits}
\slidesScale{01}{21}
}
\subsection{Steuerwerk}
\slidesScale{01}{22}
\subsection{CPU}
\slidesScale{01}{23}
\subsection{E/A-Werk und Bus}
\slidesScale{01}{24}

\section{Architekturverfeinerung 2: Einführung von Registern}
\slidesScale{01}{25}
\unimptnt{
\subsection{Von-Neumann-Rechner mit Register-Files}
\slidesScale{01}{26}
}

\section{Architekturverfeinerung 3: Erweiterung um relative Adressierung}
\slidesScale{01}{27}

\section{Universalrechner: Instruktionstypen / Befehlssatz}
\slidesScale{01}{29}

\paragraph{Beispielrechner: einfacher Von-Neumann-Beispielrechner}
\slidesScale{01}{28}

\subsection{Instruktionssatz}
\slidesScale{01}{30}

\subsection{Befehlszyklus}
\slidesScale{01}{31}

\unimptnt{
\subsubsection*{Bsp.: Interpretationszyklus des Beispielrechners}
\slidesScale{01}{32}
}

\section{Abweichungen moderner Rechner von den Von-Neumann-Prinzipien}
\slidesScale{01}{33}

\chapter{Prozessor}

\section*{Grundelemente des Prozessors} WICHTIG!
\slidesScale{02}{1}
\section{Rechenwerk}
\subsection*{Aufbau und Ansteuerung}
\slidesScale{02}{2}
\subsection*{Operationen einer (Ganzzahl-)ALU}
\slidesScale{02}{3}

\subsection{Addierer}
\subsubsection*{Addierwerk}
\slidesScale{02}{4}
\subsubsection*{Parallele Addierer: Ripple Carry Addierer (RCA)}
\slidesScale{02}{5}
\subsubsection*{Parallele Addierer: Carry-Look-Ahead-Addierer (CLA)}
\slidesScale{02}{6}

\subsection{Multiplizierer}
\unimptnt{
\subsubsection{Multiplikation von Dualzahlen}
\slidesScale{02}{7}
}

\subsubsection{Multiplizierwerk::seriell}
\slidesScale{02}{8}
\subsubsection{Multiplizierwerk::parallel}
\slidesScale{02}{9}

\unimptnt{
\subsubsection{Weitere Multiplizierwerke}
\begin{itemize}
\item Wallace Tree
\item Binary Tree
\item Booth-Algorithmus
\item Carry Save Multiplier
\end{itemize}
}

\section{Steuerwerk}

\subsection*{Grundaufgaben}
\slidesScale{02}{10}
\subsection*{Phaseneinteilung der Befehlsbearbeitung}
\slidesScale{02}{11}
\subsection*{Steuerwerk und Operationswerk}
\slidesScale{02}{12}
\subsection*{Realisierungstechniken}
\slidesScale{02}{13}

\unimptnt{
\subsection{Steuerwerk als Automat}
\subsubsection*{Finite State Machine (FSM) = endlicher Automat}
\slidesScale{02}{14}
\subsubsection*{Logikrealisierung}
\slidesScale{02}{15}
\slidesScale{02}{16}

\subsection*{Beispiele}
\begin{itemize}
\item „sequentieller Addierer“ Datenpfad\\
Zustände laut Algorithmus:
\begin{enumerate}
\item „Laden von Summand A“
\item „Laden von Summand B“
\item „Übernahme des Additionsergebnisses“
\item „Weiterschalten der Summanden um eine Stelle“
\item „Ausgabe des Ergebnisses“
\end{enumerate}
\end{itemize}
}

\subsection{Steuerwerk mit Mikroprogrammierung}

\subsubsection{Realisierung}
\slidesScale{02}{17}
\unimptnt{
\subsubsection*{Aufbau eines einfachen Mikroprogramm-Steuerwerk}
\slidesScale{02}{18}
}
\subsubsection{Mikroprogrammierung: Vor- und Nachteile}
\slidesScale{02}{19}
\subsubsection{Unterschied mikroprogrammiert/mikroprogrammierbar}
\slidesScale{02}{20}

\section{Adresswerk}
(Address Unit, Address Generation Logic)
\slidesScale{02}{21}

\unimptnt{
\subsection*{Aufbau eines einfachen Adresswerks}
\slidesScale{02}{22}
\slidesScale{02}{23}
}

\chapter{Speicherhierarchie und Cache}
%\subsection*{Speicherverwaltung: Einführung}
\slidesScale{03}{2}
\section{Speicherhierarchie}
\slidesScale{03}{3}

\subsection{Ebenen}
\slidesScale{03}{4}
\slidesScale{03}{5}

\section{Cache}
\subsection{Einführung}
\subsubsection{Motivation und Grundidee}
\slidesScale{03}{6}

\subsubsection{Umsetzung der Idee}
\slidesScale{03}{7}
\slidesScale{03}{8}

\subsubsection{Hierarchie, Übertragungsmodi}
\slidesScale{03}{9}

\subsection{Lokalität}
\subsubsection{Prinzip}
\slidesScale{03}{10}
Beispiel Lokalität: Wenn man bei C Arrays multipliziert, kann je nach Schachtelung der Schleifen die Laufzeit leiden. In C werden Arrays zeilenweise gespeichert, multipliziert man die Arrays allerdings spaltenweise, so muss viel gesprungen werden, da die Speicherblöcke weiter auseinander liegen.

\subsubsection{Zeitliche Lokalität}
(Temporal Locality)
\slidesScale{03}{11}

\subsubsection{Räumliche Lokalität}
(Spatial Locality)
\slidesScale{03}{12}

\subsubsection{Bedeutung der Lokalität}
\slidesScale{03}{13}

\subsubsection{Umsetzung des Lokalitätsprinzips}
\slidesScale{03}{14}

\subsection{Zugriff}
\subsubsection{Anordnung bezüglich CPU und Hauptspeicher}
\slidesScale{03}{15}
\subsubsection{Fehlzugriffe}
\slidesScale{03}{16}
weiterhin:
\slidesScale{03}{17}
\subsubsection*{Compulsory-Miss}
\slidesScale{03}{18}

\subsubsection{Lesezugriffe}
\slidesScale{03}{19}

\subsubsection{Schreibzugriffe}
\subsubsection*{Rückschreibestrategien}
\slidesScale{03}{20}
\subsubsection*{Durchschreibverfahren}
\slidesScale{03}{21}
\subsubsection*{Rückschreibeverfahren}
\slidesScale{03}{22}
\subsubsection*{Write Hit und Write Miss}
\slidesScale{03}{23}

\subsection{Kohärenz und Konsistenz}
\subsubsection{Kohärenz}
„Alle Speicherzugreifenden bekommen den richtigen Wert, wenn sie auf den Speicher zugreifen“
\unimptnt{
\subsubsection*{Kohärenz Beispiel mit Write-Through-Caches}
\slidesScale{03}{24}
\subsubsection*{Kohärenz Beispiel mit Write-Back-Caches}
\slidesScale{03}{25}
}

\subsubsection*{Speicherkohärenz}
\slidesScale{03}{26}

\subsubsection*{Definition Kohärenz}
\slidesScale{03}{27}
\slidesScale{03}{28}

\subsubsection{Konsistenz}
\slidesScale{03}{29}

\unimptnt{
\subsubsection*{Speicher-Konsistenz}
\slidesScale{03}{30}
}

\subsubsection{Kohärenz-Protokolle}
\slidesScale{03}{31}
\unimptnt{
\subsubsection*{Vereichnisbasiert}
\slidesScale{03}{32}

\subsubsection*{Bus-Snooping}
\slidesScale{03}{33}

Möglichkeiten:
\slidesScale{03}{34}

\subsubsection*{MESI-Kohärenzprotokoll}
\slidesScale{03}{35}
Verbesserungen:
\slidesScale{03}{36}

\subsubsection*{Update-basierte Protokolle}
\slidesScale{03}{37}
}

\section{Cacheorganisation}

\subsection{Realisierung}
\slidesScale{03}{38}

\subsection{Aufbau eines Cache-Speichers}
\slidesScale{03}{39}

\subsection{Begriffe}
\slidesScale{03}{40}

\subsection{Direktabbildender Cache}
(Direct Mapped Cache)
\slidesScale{03}{41}
Beispiel: Direkt abgebildeter Cache mit 8 Einträgen, Speicher mit $2^5$ Adressen: Die letzten Bits der Speicheradressen bestimmen die Position im Cache: $00001$, $01001$, $10001$ und $11001$ landen im Cache-Block $001$.
\unimptnt{
\slidesScale{03}{42}
}

\subsection{vollassoziativer Cache}
(Fully Associative Cache)
\slidesScale{03}{43}

\subsection{p-fach satzassoziativer Cache}
\slidesScale{03}{44}
\unimptnt{
\slidesScale{03}{45}

\subsection{Cache-Mapping: zusammenfassendes Beispiel}
\slidesScale{03}{46}
Die 2. und 3. Caches geben dem Cache also eine Struktur vor, die die Nachteile des voll-assoziativer Caches versuchen abzumildern.
}

\section{Leistungssteigerung durch Cache}

\slidesScale{03}{47}
\slidesScale{03}{48}
\unimptnt{
\slidesScale{03}{49}
}

\subsection{Gesamtgröße und Block(rahmen)größe}
\slidesScale{03}{50}
\unimptnt{
\subsubsection{Einfluss der Blockrahmengröße}
\slidesScale{03}{51}
}

\section{Nachladestrategien}
\subsection{Demand Fetching}
\slidesScale{03}{52}
\subsection{Prefetching}
\subsubsection*{Was}
\slidesScale{03}{53}
\subsubsection*{Wann}
\slidesScale{03}{54}
\unimptnt{
\slidesScale{03}{55}
}
Tagged Prefetch: Es wird nur weiter nachgeladen, wenn es sich zuvor „gelohnt“ hat ($\to$ ein Hit vorhanden war)

\section{Verdrängungsstrategien}
(Replacement Policy)
\slidesScale{03}{56}
\subsection{Random, FIFO, Round Robin}
\slidesScale{03}{57}
\subsection{LFU}
\slidesScale{03}{58}
\subsection{LRU}
\slidesScale{03}{59}
\unimptnt{
\subsubsection{Algorithmus}
\slidesScale{03}{60}
}

\chapter{Speicher und Speicherverwaltung}
\section*{Hauptspeicher}
\slidesScale{04}{1}
\section{Virtueller Speicher}
\subsection{Ziel}
\slidesScale{04}{2}
\subsection{Virtueller Speicher}
\slidesScale{04}{3}
\subsection{Speicherverwaltung: früher und heute}
\slidesScale{04}{4}
\subsubsection{Begriffe}
\slidesScale{04}{5}
\slidesScale{04}{6}
\slidesScale{04}{7}
\unimptnt{
\subsubsection{Anordnung von Cache und MMU}
\slidesScale{04}{8}
}

\section{Dynamische Zuordnungsverfahren}
\subsection{Relocation}
\slidesScale{04}{9}
\unimptnt{
\subsubsection*{Beispiel}
\slidesScale{04}{10}
}
\subsubsection{Probleme}
\slidesScale{04}{11}
\subsection{Paging}
\subsubsection{Methode}
\slidesScale{04}{12}
\subsubsection{Prinzip der virtuellen Speicherverwaltung}
\slidesScale{04}{13}
\subsubsection{Organisation der Seitentabellen}
\slidesScale{04}{14}

\subsubsection{Einstufige Adressumsetzung}
\slidesScale{04}{15}

\subsubsection*{Adressierung / Adressumsetzung}
\slidesScale{04}{16}
\unimptnt{
\slidesScale{04}{17}
\subsubsection*{Demand Paging}
„Seitennachladen nach Bedarf“: Nur die Seiten (Programmelemente) werden nachgeladen, die benötigt werden.
\subsubsection*{Demand Paging - Page Faults}
\slidesScale{04}{18}
\subsubsection*{Beispiel}
\begin{itemize}
\item 32 Bit virtuelle Adresse ($\to$ 4 GiByte virtueller Adressraum)
\item 1 MiByte ($=2^{20}$ Byte) realer Speicher
\item Seitengröße 4 KiByte ($=2^{12}$ Byte), Adressauflösung: 1 Byte\\
$\Rightarrow$ virtueller Seitennummer 32-12 = 20 Bit\\
$\Rightarrow 2^{20}$ Einträge in der Page Table
\item 1 MiByte Hauptspeicher $\Rightarrow 2^{20-12}=2^8=256$ Seitenrahmen\\
$\Rightarrow$ PTE.Rp benötigt 8 Bit = 1 Byte, 1 weiteres Byte für P, RWX etc.\\
$\Rightarrow$ $2^{20}$ Einträge a 2 Byte $\Rightarrow$ 2 MiByte für Page Table!!! (Konflikt zu 1 MiByte realer Speicher)\\
$\Rightarrow$ PT ebenfalls im virtueller Adressraum, muss auch dem Demand Paging unterliegen…
\end{itemize}
$\Rightarrow$ Dieser Aufwand wird \emph{Table Superfluity} genannt.\bigskip\\
Anderes Szenario:
\begin{itemize}
\item 32 Bit virtuelle Adresse ($\to$ 4 GiByte virtueller Adressraum)
\item 256 MiByte ($=2^{28}$ Byte) realer Speicher
\item Seitengröße 4 KiByte ($=2^{12}$ Byte), Adressauflösung: 1 Byte\\
$\Rightarrow$ virtueller Seitennummer 32-12 = 20 Bit\\
$\Rightarrow 2^{20}$ Einträge in der Page Table
\item 1 MiByte Hauptspeicher $\Rightarrow 2^{28-12}=2^{16}=65.536$ Seitenrahmen\\
$\Rightarrow$ PTE.Rp benötigt 16 Bit = 2 Byte, 1 weiteres Byte für P, RWX etc.\\
$\Rightarrow$ $2^{20}$ Einträge a 3 Byte $\Rightarrow$ 3 MiByte für Page Table\\
$\Rightarrow$ PT kann im realen Hauptspeicher gehalten werden
\end{itemize}
$\Rightarrow$ Dieser Aufwand wird \emph{Table Superfluity} genannt.
}
\subsubsection{Zweistufige Adressumsetzung}
\slidesScale{04}{19}

\subsubsection{Inverted Page Table}
\slidesScale{04}{20}
\unimptnt{
\slidesScale{04}{21}
\subsubsection*{Adressumsetzung}
\slidesScale{04}{22}
\slidesScale{04}{23}
}
\subsection{Segmentation}
\subsubsection{Methode}
\slidesScale{04}{24}
\unimptnt{
\subsubsection{Adressierung / Adressumsetzung}
\slidesScale{04}{25}
}
\subsubsection{Externe Fragmentierung}
\slidesScale{04}{26}
\unimptnt{
\subsubsection{Kompaktierung}
\slidesScale{04}{27}
}

\subsection{Paged Segmentation}
\slidesScale{04}{28}
\unimptnt{
\slidesScale{04}{29}
}

\section{Speicherverwaltungsverfahren}
\subsection{Überblick}
\slidesScale{04}{30}
\subsection{Fetch / Placement}
\slidesScale{04}{31}

\section{Translation Lookaside Buffer}
\slidesScale{04}{32}

\unimptnt{
\subsection{Adressumsetzung}
\slidesScale{04}{33}
}

\subsection{Ansatz}
\slidesScale{04}{34}

\subsection{Probleme}
\slidesScale{04}{35}
\unimptnt{
\slidesScale{04}{36}
}
$\Rightarrow$ In modernen Systemen mit stets wechselnden Prozessen lohnt sich kein zu großer TLB, da er mit jedem Prozess neu beschrieben wird. Wird also der Prozess gewechselt, was regelmäßig geschieht, ist der gesamte TLB ungültig… obwohl der vorhergehende Prozess ihn wahrscheinlich gar nicht komplett füllen konnte.

\chapter{Speicher und Speicherverwaltung}
\unimptnt{
\section{Einführendes Beispiel}
\slidesScale{05}{1}
}

\section{Prinzip des Pipelining}

\unimptnt{
\subsection{idealisiertes Prozessor-/Speichersystem}
\slidesScale{05}{2}
}
\subsubsection{Typischer Befehlszyklus}
\slidesScale{05}{3}

\subsection{Einführung}
\slidesScale{05}{4}
\slidesScale{05}{5}

\subsection{Grundidee}
\slidesScale{05}{6}

\subsection{Fließband-, Pipeline-, Ablaufdiagramm}
\slidesScale{05}{7}

\section{Leistungsfähigkeit und -parameter}
\subsection{Steigerung der Leistungsfähigkeit}
\unimptnt{
\slidesScale{05}{8}
}
\slidesScale{05}{9}
\slidesScale{05}{10}
\subsection{Pipeling-Paradoxon}
\slidesScale{05}{11}
Loop-Unrolling: anstatt bspw. eine 10-fach Schleife als solche zu notieren (while/for), diese in 10 Schritte sequentiell hin schreiben.

\section{Konflikte beim Pipelining}
\subsection{Überblick}
\slidesScale{05}{12}
\slidesScale{05}{13}

\subsection{Strukturelle Konflikte}
\subsubsection{Belegungsmatrix}
\slidesScale{05}{14}
\unimptnt{
\subsubsection*{struktureller Konflikt}
\slidesScale{05}{15}
}
\subsubsection{Entwurfsziele}
\slidesScale{05}{16}
\subsubsection{Uni- und Multifunktionale Pipelinestruktur}
\slidesScale{05}{17}
\subsubsection{Lineares Pipelineschema}
\slidesScale{05}{18}
\unimptnt{
\subsubsection{Auswirkung einer nichtlinearen, multifunktionalen Pipeline}
\slidesScale{05}{19}
}
\subsubsection{Weitere Entwurfsziele}
\slidesScale{05}{20}
\subsubsection{Ursache und Lösung}
\slidesScale{05}{24}
\unimptnt{
\subsubsection{Beispielarchitektur}
\subsubsection*{Datenpfad}
\slidesScale{05}{21}
\subsubsection*{Hardware-Belegungsmatrix}
\slidesScale{05}{22}
\subsubsection*{Ablaufdiagramm}
\slidesScale{05}{23}
\subsubsection*{Lösung}
\slidesScale{05}{25}
}

\subsubsection{Lösung Harvard-Architektur}
\slidesScale{05}{26}

\unimptnt{
\subsubsection{Leistungsfähikeitsgewinn durch Harvard-Architektur}
\subsubsection*{Beispielsituation}
\slidesScale{05}{27}
\subsubsection*{Beispielrechnung}
\slidesScale{05}{28}
}

\subsection{Datenkonflikte}
\slidesScale{05}{29}

\unimptnt{
\subsubsection*{Beispiel}
\slidesScale{05}{30}
}
\subsubsection{Lösungsansatz mittels Software}
\slidesScale{05}{31}
\subsubsection{Lösungsansatz mittels Hardware}
\slidesScale{05}{32}
\slidesScale{05}{33}
\subsubsection*{Aufwand für Forwarding}
\slidesScale{05}{34}

\subsubsection{Vergleich der Lösungsansätze}
\slidesScale{05}{35}

\subsection{Steuerkonflikte}
\slidesScale{05}{36}

\subsubsection{Auswirkungen}
\slidesScale{05}{37}
\slidesScale{05}{38}

\subsubsection{Lösungansätze mittels Software}
\slidesScale{05}{39}

\chapter{Pipelining}
\unimptnt{
\section{Einführung}
--
}

\section{Statische Vorhersage}
\subsection{Stall/Freeze}
\slidesScale{06}{2}
\subsection{Delayed Branches}
\slidesScale{06}{3}
\subsection{Predict (Not) Taken}
\slidesScale{06}{4}
\unimptnt{
\subsection{Erweiterung}
\slidesScale{06}{5}
}






\chapter{Sprungvorhersage}
\unimptnt{
\section{Einführung}
--
}

\section{Statische Vorhersage}
\subsection{Stall/Freeze}
\slides{06}{2}
\subsection{Delayed Branches}
\slides{06}{3}
\subsection{Predict (Not) Taken}
\slides{06}{4}
\unimptnt{
\subsection{Erweiterung}
\slides{06}{5}
}

\subsection{Compiler-adressiert}
\slides{06}{6}

\subsection{Diskussion}
\slides{06}{7}

\section{Dynamische Vorhersage}
\slides{06}{8}

\subsection{1-Bit-Prädikator}
\slides{06}{9}
\subsubsection{Vor-/Nachteile}
\slides{06}{10}

\subsection{2-Bit-Prädikator}
\slides{06}{11}
\subsubsection{Varianten}
\slides{06}{12}
\subsection{Vorteile}
\slides{06}{13}
\subsubsection{Probleme}
\slides{06}{14}

\unimptnt{
\subsection{n-Bit-Prädikator}
\slides{06}{15}
\subsubsection{Weiter Prädikatoren}
\slides{06}{16}
}

\subsection{Branch History Table / Branch Target Buffer}
\subsubsection{Branch History Table - BHT}
\slides{06}{17}
Puffer wichtig, weil die Vorhersage sofort erreichbar sein muss (nicht im Hauptspeicher, der ggf. auch virtueller Speicher und damit langsam sein kann).
\subsubsection{Branch Target Buffer - BTB}
\slides{06}{18}
\subsubsection*{Vorteile/Probleme}
\slides{06}{19}
\unimptnt{
\subsubsection*{Variationsmöglichkeiten}
\slides{06}{20}
}

%\subsection{Leistungsfähigkeit}

\chapter{Superskalare Architekturen, Parallelisierung}
\section{Einführung}
\subsection{Erhöhung der Rechenleistung: Ansätze}
\slides{07}{1}

\subsection{Superskalar Definition}
\slides{07}{2}
\subsection{Parallele Strukturen bei Einprozessorsystemen}
\slides{07}{3}

\section{Superpipeline-Architektur (Superpipelining)}
\slides{07}{4}

\subsection{Diskussion}
\slides{07}{5}

\section{VLIW-Architektur}
\slides{07}{6}
\slides{07}{7}

\unimptnt{
\subsection{EPIC-Mikroarchitekturen}
\slides{07}{8}
}

\subsection{Probleme}
\slides{07}{9}

\section{Superskalar-Architektur}

\subsection{Architektur}
\slides{07}{10}

\subsubsection*{Gemeinsame Pipeline}
\slides{07}{11}

\subsubsection*{Superskalare Prozessor-Pipeline}
\slides{07}{12}
\unimptnt{
\slides{07}{13}
}

\subsubsection{Superskalarität: Scheduler}
\slides{07}{14}

\unimptnt{
\subsubsection{Superskalar-Architektur}
\slides{07}{15}
}

\subsubsection[Superpipeline <-> Supersklare Architektur]{Superpipeline $\leftrightarrow$ Supersklare Architektur}
\slides{07}{16}

\subsubsection{Superpipeline \& Superskalar kombiniert}
\slides{07}{17}

\subsubsection{Zusammenfassung}
\slides{07}{18}

\subsection{Befehlszuteilung}
\slides{07}{19}

\subsubsection{statisch (In-Order-Execution)}
\slides{07}{20}
\slides{07}{21}

\subsubsection{dynamisch (Out-Of-Order-Execution)}
\slides{07}{22}
\slides{07}{23}
\unimptnt{
\subsubsection*{Teilaufgaben}
\slides{07}{24}
\slides{07}{25}
}

\subsubsection{dynamisch: Scoreboard}
\slides{07}{26}

\unimptnt{
\subsubsection*{Ablaufüberwachung}
\slides{07}{27}
}

\subsubsection{dynamisch: Tomasolu-Methode}
\slides{07}{28}

\chapter{Leistungsbewertung}
\section{Einführung}
\subsection{Trends}
\slides{08}{2}

\unimptnt{
\subsection{stete Leistungszunahme}
\slides{08}{3}

\subsection{Bewertung der Leistungsfähigkeit}
\slides{08}{4}
\slides{08}{5}
}

\section{Einfache Metriken: MIPS, FLOPS und Co.}

\slides{08}{6}
\slides{08}{7}

\subsection{Einfache Hardwaremaßzahlen}

\subsubsection{MIPS}
\slides{08}{8}
\unimptnt{
\subsubsection{FLOPS, OPS, TOPS}
\slides{08}{9}
}
\subsubsection{MPS, FLOPS \& Co. in der Kritik}
\unimptnt{
\slides{08}{10}
}
\slides{08}{11}

\section{Auswertung von Hardwaremaßnahmen}

\subsection{Mittlere Operationszeit}
\slides{08}{12}

\subsection{Ausführungszeit (Execution Time)}
\unimptnt{
\slides{08}{13}
}
\slides{08}{14}

\subsection{Parallelitätsprofil}
\slides{08}{15}

\subsection{Einfache Hardwaremaßzahlen:  CPU-Zeit und CPI}
\slides{08}{16}
\slides{08}{17}

\subsection{Befehlsverteilung und der CPI-Wert}
\slides{08}{18}

\subsection{Programmabhängiges Leistungsmodell}
\slides{08}{19}

\subsection{Minimierung der Anzahl ausgeführter Befehle}
\slides{08}{20}
\slides{08}{21}

\subsection{Minimierung der Taktzyklen pro Durchschnittsbefehl}
\slides{08}{22}

\subsection{Minimierung der Taktzykluszeit}
\slides{08}{23}

\subsection{(Normierte) Maximalleistung}
\slides{08}{24}
\slides{08}{25}

\subsection{Prozessor-Speicher-Leistungsbewertung}
\slides{08}{26}

\subsection{Zusammenfassung: Auswertung von Hardwaremaßnahmen}
\slides{08}{27}

\section{Parallelverarbeitung, Amdahl, Gustafson}
\subsection{Speedup}
\slides{08}{28}

\subsection{Effizienz}
\slides{08}{29}

\subsection{Ahmallsches Gesetz bei Parallelverarbeitung}
\slides{08}{30}
\unimptnt{
\slides{08}{31}
}
\slides{08}{32}
\subsubsection{Ahmdalsches Gesetz}
\slides{08}{33}

\subsection{Gustafsons Gesetz}
\subsubsection{Ansatz}
\slides{08}{34}
\subsubsection{Erkenntnisse}
\slides{08}{35}
\unimptnt{
\subsubsection{Herleitung}
\slides{08}{36}
}
\subsection{Ahmdal vs Gustafson}
\unimptnt{
\slides{08}{37}
}
Hauptunterschied Ahmdal vs Gustafson:\\
Ahmdal schließt von der Leistung eines seriellen Systems auf die eines parallelen Systems.\\
Gustafson schließt von der Leistung eines parallelen Systems, wie langsam es auf einem seriellen System laufen würde.\\
Dem entsprechend kommen unterschiedlich Annahmen heraus.

\subsection{Skalierbarkeit eines Parallelsystems}
\slides{08}{38}

\section{Benchmarks}
\subsection{Benchmark-Programme}
\slides{08}{39}
\unimptnt{
\subsection{Benchmark-Typen}
\subsubsection{Reale Applikationen}
\slides{08}{40}
\subsubsection{Toy Benchmarks}
\slides{08}{41}
\subsubsection{Kernels}
\slides{08}{42}
\slides{08}{43}
\subsubsection{Synthetische Benchmarks}
\slides{08}{44}
\subsubsection*{Whetstone / 1976}
\slides{08}{45}
\subsubsection*{Dhrystone / 1984}
\slides{08}{46}
}
\subsection{Standardisierte Benchmarks}
\slides{08}{47}
\unimptnt{
\slides{08}{48}
\subsubsection{SPEC}
\slides{08}{49}
\subsubsection*{CPU2006}
\slides{08}{50}
}

\section{Abschließendes}
\unimptnt{
\subsection{Benchmark-Wahn und seine Probleme}
\slides{08}{51}
\subsection{Leistungsmessung: übliche Trugschlüsse}
\slides{08}{52}
}
\subsection{Leistungsmessung: Zusammenfassung}
\slides{08}{53}














\chapter{Informationstransport: Bus, Protokolle}
\section{Bus: Einführung}
\subsection{Definition}
\slides{09}{2}
\subsection{Generelle Struktur}
\slides{09}{3}
\subsection*{Was definiert einen Bus?}
\slides{09}{4}
\subsection{Vor- und Nachteile}
\slides{09}{5}
\subsection{Terminologie}
\slides{09}{6}
\slides{09}{7}
\slides{09}{8}
\slides{09}{9}
\subsection{Physikalische Struktur: parallel und seriell}
\slides{09}{10}

\section{Hierarchie}
\subsection*{Hierarchische Organisation}
\slides{09}{11}
\subsection{Prozessorbus}
\slides{09}{12}
\subsection{Speicherbus}
\slides{09}{13}
\subsection{Peripheriebus}
\slides{09}{14}
\subsection{Ein-/Ausgabebus}
\slides{09}{15}

\unimptnt{
\subsection{Alternative Einteilung}
\slides{09}{16}
}

\section{Topologie und Kopplung}
\subsection{Topologien}
\slides{09}{17}
\subsubsection{Repeater}
\slides{09}{18}
\subsubsection{Hub (Nabe/Mittelpunkt)}
\slides{09}{19}
\subsubsection{Bridge (Brücke)}
\slides{09}{20}

\section{Dedizierung}
\subsection{Dedizierter Bus}
\slides{09}{21}
\subsection{Nichtdedizierter Bus}
\slides{09}{22}

\section{Partitionierung}
\subsection{Ressourcenpartitionierter Bus}
\slides{09}{23}
\unimptnt{
\subsubsection*{Charakteristika}
\slides{09}{24}
}
\subsection{Funktionspartitionierter Bus}
\slides{09}{25}

\section{Transaktion und Übertragungsarten}
\subsection{Überblick}
\slides{09}{26}
\subsubsection*{(Nicht-)multiplexiert}
\slides{09}{27}
\unimptnt{
\subsection{Multiplexbus}
\slides{09}{28}
}




\chapter{E/A und Interrupts}
\section{Einführung}
\slides{10}{2}
\slides{10}{3}

\section{Programmierte E/A}
\slides{10}{4}
\subsection{Memory Mapped I/O}
\slides{10}{5}
\subsection{I/O Mapping (Isolierte Adressierung)}
\slides{10}{6}
Weiterer Vorteil: Adressräume können dynamisch verschoben werden.
\unimptnt{
\subsection{Befehle}
\slides{10}{7}
}
\subsection[E/A Testschleife -> Polling (Busy Waiting)]{E/A Testschleife $\to$ Polling (Busy Waiting)}
\slides{10}{8}
\slides{10}{9}

\section{Interrupts}
\slides{10}{10}
\subsection{Grundmethode (technische Lösung)}
\slides{10}{11}
\unimptnt{
\subsection{Sprung zur ISR}
\slides{10}{12}
}
\subsection{Identifikation der Quelle}
\slides{10}{13}

\subsubsection{Polling}
\slides{10}{14}
\unimptnt{
\slides{10}{15}
Es wird rundum abgefragt, wer den Interrupt ausgelöst hat. Wenn einer gefunden wurde, wird an der Stelle weiter gesucht.
\slides{10}{16}
Es wird rundum abgefragt, wer den Interrupt ausgelöst hat. Wenn einer gefunden wurde, wird wieder von vorne rundum abgefragt (die letzten können ggf. nie erreicht werden).
}

\subsubsection{Daisy Chain}
\slides{10}{17}

\subsubsection{Multilevel Interrupt}
\slides{10}{18}

\subsection{Ermittlung der Adresse der ISR}
\slides{10}{19}
\subsubsection{Berechnung der ISR-Startadresse}
\slides{10}{21}

\subsection{(Ausnahme-)Vektortabelle}
\slides{10}{20}


\subsection{Maskierung}
\unimptnt{
\slides{10}{22}
}
\slides{10}{23}

\unimptnt{
\subsection{Programmierbare Interrupt-Controller (PIC)}
\slides{10}{24}
}

\subsection{Programmierte E/A mit Interrupts: Vor-/Nachteile}
\slides{10}{25}

\unimptnt{
\subsection{Noch mehr Verbesserungen?}
\slides{10}{26}
}

\section{Direct Memory Access (DMA)}
\slides{10}{27}
\subsection{Definition}
\slides{10}{28}

\subsection{DMA-Controller (DMAC)}
\slides{10}{29}

\subsection{Konfliktlösung bei Speicherzugriff}
\slides{10}{30}

\section{E/A-Prozessoren}
\subsection{Prinzip}
\slides{10}{31}

\unimptnt{
\subsection{Device Driver}
\slides{10}{32}
}

\subsection[Anschlussschemata IOP <=> Geräte]{Anschlussschemata IOP $\Leftrightarrow$ Geräte}
\slides{10}{33}

\renewcommand{\unimptnt}[1]{}
\part{Praktikum}
\setCustomSlidePath{Praktikum/RA_043_P00D}

\chapter{Einführung}
\section{Mikrocontroller}
\slideScale{2}
\subsection{Unterscheidung und Architektur}
\slideScale{3}

\unimptnt{
\subsection{Taktung}
\slideScale{4}
\subsection{Typische Peripherie}
\slideScale{5}
}
\subsection{Einsatzbeispiele}
\slideScale{6}
\subsection*{Computerklassen/-typen}
\begin{itemize}
\item Wegwerfcomputer
\item Microcontroller, embedded Computer
\item Mobile Computer
\item Personalcomputer
\item Server
\item Großrechner
\end{itemize}

\subsection{Speicherbedarf}
\slideScale{7}
\slideScale{8}

\unimptnt{
\chapter{Fallbeispiel Atmel AVR}
\section{AVR und AVR-Architektur im Überblick}
\subsection{Allgemeines}
\slideScale{9}
\subsection{CPU und Archtitektur}
\slideScale{10}
\slideScale{11}
\slideScale{12}
\subsection{CPU zu Statusregister SREG}
\slideScale{13}
\subsection{Interne Perepherie}
\slideScale{14}

\section{Speichertypen}
\subsection{Speicher}
\slideScale{15}
\slideScale{16}
\subsection{Stack}
\slideScale{17}

\section{Beispielcontroller}
\subsection{Typenbeispiele}
\slideScale{18}
\subsection{Blockschaltbild und Besonderheiten}
\slideScale{19}
\subsection{Gehäuse und Anschlüsse}
\slideScale{20}

\section{Entwicklungswerkzeuge}
\subsection{Demoplattform AVR Butterfly}
\slideScale{21}
\subsection{Hardware}
\slideScale{22}
\subsubsection{Experimentalplattform STK500/501}
\slideScale{23}
\slideScale{24}
\subsection{Software}
\slideScale{25}

\section{Baugruppen und deren Nutzung und Programmierung}
\subsection{System- und Takt-Anschlüsse}
\slideScale{26}
\subsection{Bidirektionale Ports}
\slideScale{27}
\slideScale{28}
\slideScale{29}
\slideScale{30}
uint8\_t für 8-Bit unsigned Integer - gut, da AVR 8-Bit ist.
\slideScale{31}
\slideScale{32}
\slideScale{33}
}
Wichtig: Bit-Tricks !!\\
1<<3 : 3. Bit wird auf 1 gesetzt; |= (bitweises Oder) belässt alle Bits gleich, bis auf das 3. (in dem Bsp.), das wird 1.\\
\textasciitilde() (bitweise Negation); \&= belässt alle Bits gleich, bis auf das 3., das wird 1.\\
\^{}= (exklusives Oder) fürs Invertieren
\unimptnt{
\slideScale{34}
\slideScale{35}
\slideScale{36}
\slideScale{37}
\slideScale{38}

\subsection{Alternative Nutzung der Universalports}
\slideScale{39}
\slideScale{40}
\slideScale{41}
\slideScale{42}

\subsection{Interrupts}
\subsubsection{Grundprinzip}
\slideScale{43}

\subsubsection{Interruptquellen}
\slideScale{44}
\subsubsection*{Beispiel ATmega128}
\slideScale{45}
\slideScale{46}
\subsubsection{Anforderungen an Interrupt-Routinen}
\slideScale{47}
Interrupt in C: mit globalen Variablen
\subsubsection{Beispiel für Interruptnutzung}
\slideScale{48}
Hinweis Atmel-Board: 0 ist "1" und 1 ist "0" (im Sinne, dass bei 0 die Lampe leuchtet)
\slideScale{49}
\slideScale{50}
\slideScale{51}

\subsection{Timer/Counter}
\slideScale{52}
\subsubsection{Prescaler (Vorteiler)}
\slideScale{53}
\subsubsection{Betriebsmodi}
\slideScale{54}

\subsubsection{Normal Mode}
\slideScale{55}
\subsubsection{CTC-Mode}
\slideScale{56}
\subsubsection{PWM-Modus}
\slideScale{57}
\subsubsection{Beispiel 1 für Timer/Counter-Nutzung}
\slideScale{58}
\subsubsection{Beispiel 2 für Timer/Counter-Nutzung}
\slideScale{59}
\slideScale{60}
\slideScale{61}
\slideScale{62}
\slideScale{63}
}

%\newpage
%\printbibliography

\end{document}
